% ============================================================================
% preface.tex — Preface
% ============================================================================

\chapter*{Preface}
\addcontentsline{toc}{chapter}{Preface}
\markboth{Preface}{Preface}

\begin{center}
\begin{tikzpicture}
  \fill[ModuleBlue!10, rounded corners=10pt] (-7,-0.6) rectangle (7,0.6);
  \node[font=\Large\bfseries\sffamily, text=NavyBlue] 
    {Welcome to Your Computer Science Adventure!};
\end{tikzpicture}
\end{center}

\vspace{8pt}

Hello there, future tech expert! \faRocket

Welcome to \textbf{Foundations of Computer Science} --- your personal guide to understanding the amazing world of computers. Whether you've been using computers for years or you're just getting started, this book is designed to take you on an exciting journey from the very basics all the way to creating your own mini-projects. Think of it as your friendly companion that explains everything in plain, simple language --- no confusing jargon, we promise!

\vspace{8pt}

This book is organised into \textbf{eight modules}, and each one builds on what you learned before. You'll start by discovering what a computer actually \emph{is} (spoiler: it's not magic, even though it sometimes feels like it!). Then you'll explore Windows, master keyboard shortcuts that will make you faster than your friends, create professional documents in Microsoft Word, crunch numbers in Excel, build stunning presentations in PowerPoint, and learn to navigate the internet safely. By the end, you'll bring everything together in a capstone project that shows off all your new skills.

\vspace{8pt}

Every chapter follows a simple pattern: first, we explain the \textbf{theory} --- the ``what'' and ``why.'' Then, we give you \textbf{shortcuts and tips} to work smarter, not harder. Finally, you get \textbf{hands-on lab practicals} where you actually \emph{do} things on a real computer. Oh, and keep an eye out for \textbf{Byte} --- our friendly robot mascot who pops up throughout the book with tips, fun facts, and the occasional warning. Byte's got your back!

\vspace{8pt}

One last thing: \textbf{don't worry about getting everything right the first time.} Seriously. Experiment, make mistakes, click the wrong button, accidentally change the font to something ridiculous --- and then figure out how to fix it. That's exactly how real computer scientists work. Every expert was once a beginner, and every great discovery started with someone saying, ``I wonder what happens if I try this\ldots''

\vspace{12pt}

\begin{flushright}
\textit{Happy learning!}\\[4pt]
\textbf{The Computer Science Department}
\end{flushright}

\bytetip[fun]{Hi! I'm \textbf{Byte}, your robot guide! Whenever you see me, I'll share a cool tip, a fun fact, or a helpful warning. Let's have fun learning together!}
